% Generated draft: requires mathematical and source review.
\documentclass[11pt]{article}
\usepackage[margin=25mm]{geometry}
\usepackage{fontspec}
\setmainfont{Latin Modern Roman}
\usepackage{amsmath,amssymb,mathtools,unicode-math}
\setmathfont{Latin Modern Math}
\usepackage{enumitem,needspace,xurl,hyperref}
\setlength{\parindent}{0pt}
\setlength{\parskip}{5pt}
\setlength{\emergencystretch}{3em}
\newcommand{\N}{\mathbb{N}}
\newcommand{\Z}{\mathbb{Z}}
\newcommand{\Q}{\mathbb{Q}}
\newcommand{\R}{\mathbb{R}}
\newcommand{\C}{\mathbb{C}}
\newcommand{\F}{\mathbb{F}}
\newcommand{\A}{\mathbb{A}}
\newcommand{\PP}{\mathbb{P}}
\newcommand{\E}{\mathbb{E}}
\newcommand{\eps}{\varepsilon}
\newcommand{\dd}{\,\mathrm{d}}
\begin{document}
\section*{UnsolvedMath ID: 2302057}
\textbf{Source problem code:} catalogue entry
\subsection{Definitions and mathematical statement}
For an entire function \(f\), put \(M(r,f)=\max_{|z|=r}|f(z)|\). For each real \(\alpha>1\), the source defines \(I_\alpha\) by the two simultaneous conditions
\[
\log M(r,f)=O((\log r)^{1+\alpha}),\qquad
(\log r)^\alpha=o(\log M(r,f))\quad(r\to\infty).
\]
Equivalently, the first has a constant upper bound for the displayed normalized growth eventually, and the second requires \((\log r)^\alpha/\log M(r,f)\to0\).

The source credits Hayman with the following result under the smaller-growth hypothesis \(\log M(r,f)=O((\log r)^2)\):
\[
\frac{\log|f(z)|}{\log M(|z|,f)}\longrightarrow1
\quad\text{as }|z|\to\infty\text{ outside }E,
\]
where \(E\) consists of exceptional ``circles'' whose subtended angles at \(0\) have finite sum. In particular, it records the same ratio limit for \(z=re^{i\theta}\), \(r\to\infty\), for Lebesgue-almost every \(\theta\). The first question is whether these conclusions continue to hold for every \(f\in I_\alpha\), for \(\alpha>1\).

For meromorphic \(f\), let \(T(r,f)\) be the Nevanlinna characteristic and let
\[
\delta(a,f)=1-\limsup_{r\to\infty}\frac{N(r,a)}{T(r,f)}
\quad(a\in\widehat\C),
\]
where \(N(r,a)\) is the integrated counting function with multiplicity, counting poles for \(a=\infty\). The source credits Anderson--Clunie with: if \(T(r,f)=O((\log r)^2)\), there is at most one deficient value, and every \(a\) with \(\delta(a,f)>0\) is asymptotic and satisfies
\[
f(re^{i\theta})\longrightarrow a\quad(r\to\infty)
\quad\text{for almost every }\theta.
\]
Limits to \(\infty\) are understood spherically. An asymptotic value is a limit along some curve tending to infinity. The second question asks for the corresponding extension of these conclusions. A natural meromorphic analogue of \(I_\alpha\) would require
\[
T(r,f)=O((\log r)^{1+\alpha}),\qquad
(\log r)^\alpha=o(T(r,f)),
\]
but the source explicitly defines \(I_\alpha\) only for entire functions, so this substitution is a conditional interpretation rather than an unqualified source definition.

J. M. Anderson asks whether the results depend on small growth or instead on growth smoothness. The source calls Hayman's functions \(I_1\), although its displayed definition would also impose a lower-growth condition; the precise intended inclusion requires checking. The exceptional ``circles'' also need a disk/covering convention for the stronger off-exceptional-set statement. The supplied OPEN-TRIAGE summary reports open status in the 2018 edition and no recent resolution located.
\subsection{Short English statement}
Do the maximum-modulus and deficient-value conclusions known at growth \(O((\log r)^2)\) extend to the two-sided growth classes \(I_\alpha\), \(\alpha>1\)? The meromorphic analogue and exceptional-circle convention need clarification.
\subsection{Sources}
\begin{itemize}
\item Local UnsolvedMath registry, permanent ID \(2302057\), exact source code \texttt{catalogue entry}, ``Research Problems in Function Theory --- Problem 2.57.''
\item Hayman, \emph{Research Problems in Function Theory} (2018), Problem 2.57, attributed to J. M. Anderson; \url{https://arxiv.org/abs/1809.07200}. Source-TeX extraction; recorded access 2026-07-29.
\item The statement attributes the cited small-growth results to Hayman and Anderson--Clunie without complete references. The supplied OPEN-TRIAGE summary gives no separate current-status citation.
\end{itemize}
\end{document}
